The accurate estimation of battery state of health (SOH) is a prerequisite
for the safe and efficient operation of lithium-ion batteries in electric
vehicles (EVs), grid-scale energy storage, and consumer
electronics~\cite{lipu2018review}.
SOH is the ratio of a cell's present usable capacity to its original rated
capacity; as a cell is charged and discharged over hundreds of cycles this
ratio falls, and the point at which it drops too far sets the battery's usable
lifetime and dictates when it must be replaced~\cite{birkl2017degradation}.
Predicting how SOH will decline, ideally early and for each individual cell, is
therefore central to managing cost, safety, and warranty risk in every large
battery deployment.

Over the past decade, deep learning has produced increasingly capable SOH
prediction models, progressing from standard long short-term memory (LSTM)
networks to attention-augmented variants~\cite{qu2019neural} and models whose
hyperparameters are tuned by metaheuristic
search~\cite{ponnambalam2024battery,ponnambalam2026delstm}.
Yet almost all of these methods are developed and benchmarked in the same
narrow setting: small cylindrical 18650 cells of 1 to 3\,Ah (LCO, NMC, or NCA
chemistry) cycled at constant current in the laboratory.
The NASA Prognostics Repository~\cite{saha2007battery}, and in particular its
cells B05, B06, B07, and B18, has been the field's de facto benchmark since
2007, appearing in roughly 30.8\% of all published SOH and remaining-useful-life
(RUL) studies~\cite{xu2025review,bairwa2025cycle}.

This near-exclusive reliance on small laboratory cells creates a gap that only
becomes visible at deployment.
The batteries that actually run electric vehicles and grid storage are
large-format prismatic or pouch cells of 10 to 200\,Ah~\cite{forgez2010thermal},
operated across a wide range of temperatures and charge and discharge rates.
Crucially, cells that are nominally identical, built to the same design on the
same production line, can age very differently because of small, unavoidable
variations introduced during manufacturing~\cite{birkl2017degradation}.
This cell-to-cell variability, of which manufacturing spread is a likely
contributor, is largely absent from the hand-picked, smoothly degrading cells
that populate standard benchmarks, yet it is precisely what a
manufacturer or fleet operator must contend with: warranties, safety margins,
and replacement schedules all depend on predicting the life of each individual
cell, not an average one.
Just how severe this variability is on large-format cells, and whether it,
rather than the choice of algorithm, sets the ultimate limit on prediction,
remains unknown.

A second, more subtle problem lies in how the field measures success.
Almost all of this literature ranks methods by one-step-ahead error: given a
cell's observed history, predict its SOH one cycle into the future.
On the smooth, slowly changing trajectories of large-format cells, however, this
is barely harder than assuming that the next cycle simply repeats the last,
which raises a basic question: can a metric so close to a trivial baseline
actually tell one method apart from another?
If it cannot, then a leaderboard built on it is uninformative, and predictors
need to be judged by something that matters in the field instead, the practical
capabilities they enable: predicting a brand-new cell before it has aged,
attaching trustworthy confidence bounds to each forecast, flagging defective
cells early, and forecasting remaining useful life.

The objective of this paper is to close this gap between benchmark practice and
deployment reality.
Using an 18-month dataset of 20 large-format LiFePO$_4$ prismatic cells
(50\,Ah) cycled under three temperature and C-rate
conditions~\cite{ponnambalam2026sitdata}, we set out to establish how far
cell identity, rather than model choice, governs prediction on
industrial-format cells, and to determine what a prediction method must provide
to be useful once these cells are deployed.
In place of a single accuracy score, we take as our goal a model that can
predict an unseen cell from only its earliest cycles, report and correct its own
uncertainty, remain reliable when the available training data come from a
different battery chemistry, and extend to the related deployment tasks of early
defect screening and remaining-useful-life forecasting.

The remainder of this paper is organised as follows.
Section~\ref{sec:literature} reviews the related literature.
Section~\ref{sec:gap} states the research gaps and summarises the contributions
of this work.
Section~\ref{sec:methodology} describes the datasets, the DE-LSTM baseline, and
the evaluation protocol.
Section~\ref{sec:problem} characterises the deployment challenges empirically on
the SIT cells.
Section~\ref{sec:method} presents the proposed hybrid model.
Section~\ref{sec:results} reports its evaluation, including the ablation study and
the cross-dataset validation.
Section~\ref{sec:screening} shows how the manufacturing-defective cells can be
detected early from the discharge voltage curve.
Section~\ref{sec:rul} extends the hybrid to remaining-useful-life forecasting.
Section~\ref{sec:discussion} discusses implications for deployment, and
Section~\ref{sec:conclusion} concludes.
